% =============================================================================
% Anhang
% -----------------------------------------------------------------------------
% Hierhin gehören u.a.:
%   - KI-Prompts (Dokumentationspflicht bei KI-Nutzung) -- in Kurzform (WI)
%     oder als vollstaendiger Verlauf, wenn deine Pruefenden das verlangen.
%     Manche kuendigen bei fehlender Dokumentation "nicht bestanden" an.
%   - eigene Vorarbeiten (Expose, Vorstudie) benennen -- die Eigenstaendig-
%     keitserklaerung erfasst auch Uebernahmen aus eigenen Arbeiten
%   - ergänzende Tabellen und Abbildungen
%   - Interview-Transkripte, Fragebögen (bei empirischen Arbeiten)
% =============================================================================
\section*{Anhang}
\addcontentsline{toc}{section}{Anhang}

\subsection*{A. Dokumentation der KI-Prompts}

Im Folgenden werden die im Arbeitsprozess tatsächlich verwendeten Prompts in
Kurzform dokumentiert:

\begin{enumerate}
  \item \textbf{[Werkzeug]:} \enquote{[Kurzfassung des Prompts, z.\,B. Auftrag
        zur Strukturplanung der Gliederung]}
  \item \textbf{[Werkzeug]:} \enquote{[Kurzfassung des Prompts, z.\,B. Auftrag
        zur sprachlichen Überarbeitung eines Absatzes ohne inhaltliche
        Ergänzung]}
\end{enumerate}

% Längere Prompts kannst du als Codeblock einbinden:
% \begin{lstlisting}[breaklines=true,basicstyle=\ttfamily\scriptsize]
% Hier der wörtliche Prompt ...
% \end{lstlisting}
